\section{Intersecting with an invertible sheaf and rational equivalence}
\label{section-commutativity}

\noindent
Applying the key lemma we obtain the fundamental properties of intersecting
with invertible sheaves. In particular, we will see that
$c_1(\mathcal{L}) \cap -$ factors through rational equivalence and
that these operations for different invertible sheaves commute.

\begin{lemma}
\label{lemma-commutativity-on-integral}
Let $(S, \delta)$ be as in Situation \ref{situation-setup}.
Let $X$ be locally of finite type over $S$.
Assume $X$ integral and $\dim_\delta(X) = n$.
Let $\mathcal{L}$, $\mathcal{N}$ be invertible on $X$.
Choose a nonzero meromorphic section $s$ of $\mathcal{L}$
and a nonzero meromorphic section $t$ of $\mathcal{N}$.
Set $\alpha = \text{div}_\mathcal{L}(s)$ and
$\beta = \text{div}_\mathcal{N}(t)$.
Then
$$
c_1(\mathcal{N}) \cap \alpha
=
c_1(\mathcal{L}) \cap \beta
$$
in $\CH_{n - 2}(X)$.
\end{lemma}

\begin{proof}
Immediate from the key Lemma \ref{lemma-key-formula}
and the discussion preceding it.
\end{proof}

\begin{lemma}
\label{lemma-factors}
Let $(S, \delta)$ be as in Situation \ref{situation-setup}.
Let $X$ be locally of finite type over $S$.
Let $\mathcal{L}$ be invertible on $X$.
The operation $\alpha \mapsto c_1(\mathcal{L}) \cap \alpha$
factors through rational equivalence to give an operation
$$
c_1(\mathcal{L}) \cap - : \CH_{k + 1}(X) \to \CH_k(X)
$$
\end{lemma}

\begin{proof}
Let $\alpha \in Z_{k + 1}(X)$, and $\alpha \sim_{rat} 0$.
We have to show that $c_1(\mathcal{L}) \cap \alpha$
as defined in Definition \ref{definition-cap-c1} is zero.
By Definition \ref{definition-rational-equivalence} there
exists a locally finite family $\{W_j\}$ of integral closed
subschemes with $\dim_\delta(W_j) = k + 2$ and rational functions
$f_j \in R(W_j)^*$ such that
$$
\alpha = \sum (i_j)_*\text{div}_{W_j}(f_j)
$$
Note that $p : \coprod W_j \to X$ is a proper morphism,
and hence $\alpha = p_*\alpha'$ where $\alpha' \in Z_{k + 1}(\coprod W_j)$
is the sum of the principal divisors $\text{div}_{W_j}(f_j)$.
By Lemma \ref{lemma-pushforward-cap-c1} we have
$c_1(\mathcal{L}) \cap \alpha = p_*(c_1(p^*\mathcal{L}) \cap \alpha')$.
Hence it suffices to show that each
$c_1(\mathcal{L}|_{W_j}) \cap \text{div}_{W_j}(f_j)$ is zero.
In other words we may assume that $X$ is integral and
$\alpha = \text{div}_X(f)$ for some $f \in R(X)^*$.

\medskip\noindent
Assume $X$ is integral and $\alpha = \text{div}_X(f)$ for some $f \in R(X)^*$.
We can think of $f$ as a regular meromorphic section of the invertible
sheaf $\mathcal{N} = \mathcal{O}_X$. Choose a meromorphic section
$s$ of $\mathcal{L}$ and denote $\beta = \text{div}_\mathcal{L}(s)$.
By Lemma \ref{lemma-commutativity-on-integral}
we conclude that
$$
c_1(\mathcal{L}) \cap \alpha = c_1(\mathcal{O}_X) \cap \beta.
$$
However, by Lemma \ref{lemma-c1-cap-additive} we see that the right hand side
is zero in $\CH_k(X)$ as desired.
\end{proof}

\noindent
Let $(S, \delta)$ be as in Situation \ref{situation-setup}.
Let $X$ be locally of finite type over $S$.
Let $\mathcal{L}$ be invertible on $X$.
We will denote
$$
c_1(\mathcal{L}) \cap - : \CH_{k + 1}(X) \to \CH_k(X)
$$
the operation $c_1(\mathcal{L}) \cap - $. This makes sense by
Lemma \ref{lemma-factors}. We will denote $c_1(\mathcal{L})^s \cap -$
the $s$-fold iterate of this operation for all $s \geq 0$.

\begin{lemma}
\label{lemma-cap-commutative}
Let $(S, \delta)$ be as in Situation \ref{situation-setup}.
Let $X$ be locally of finite type over $S$.
Let $\mathcal{L}$, $\mathcal{N}$ be invertible on $X$.
For any $\alpha \in \CH_{k + 2}(X)$ we have
$$
c_1(\mathcal{L}) \cap c_1(\mathcal{N}) \cap \alpha
=
c_1(\mathcal{N}) \cap c_1(\mathcal{L}) \cap \alpha
$$
as elements of $\CH_k(X)$.
\end{lemma}

\begin{proof}
Write $\alpha = \sum m_j[Z_j]$ for some locally finite
collection of integral closed subschemes $Z_j \subset X$
with $\dim_\delta(Z_j) = k + 2$.
Consider the proper morphism $p : \coprod Z_j \to X$.
Set $\alpha' = \sum m_j[Z_j]$ as a $(k + 2)$-cycle on
$\coprod Z_j$. By several applications of
Lemma \ref{lemma-pushforward-cap-c1} we see that
$c_1(\mathcal{L}) \cap c_1(\mathcal{N}) \cap \alpha
= p_*(c_1(p^*\mathcal{L}) \cap c_1(p^*\mathcal{N}) \cap \alpha')$
and
$c_1(\mathcal{N}) \cap c_1(\mathcal{L}) \cap \alpha
= p_*(c_1(p^*\mathcal{N}) \cap c_1(p^*\mathcal{L}) \cap \alpha')$.
Hence it suffices to prove the formula in case $X$ is integral
and $\alpha = [X]$. In this case the result follows
from Lemma \ref{lemma-commutativity-on-integral} and the definitions.
\end{proof}






